\documentclass{article}
\usepackage{amsmath,amssymb,amsthm}
\usepackage{algorithm}
\usepackage{algpseudocode}
\usepackage{enumitem}
\usepackage{hyperref}
\usepackage{geometry}
\geometry{margin=1in}
\newtheorem{theorem}{Theorem}
\newtheorem{lemma}{Lemma}
\newtheorem{assumption}{Assumption}
\newcommand{\E}{\mathbb{E}}
\newcommand{\R}{\mathbb{R}}

\begin{document}

\section*{Randomized Block Coordinate Descent (RBCD)}

\section*{Mathematical Formulation}
Let $f:\R^{d}\rightarrow\R$ be a convex differentiable function.  
Partition the coordinate index set $\{1,\dots ,d\}$ into $B$ disjoint blocks
\[
\{1,\dots ,d\}= \bigcup_{b=1}^{B} \mathcal{I}_{b},\qquad 
\mathcal{I}_{b}\cap \mathcal{I}_{b'}=\emptyset\;(b\neq b').
\]
Denote by $w^{(b)}\in\R^{|\mathcal{I}_{b}|}$ the subvector of $w$ that belongs to block $b$, and by
$\nabla_{b}f(w)\in\R^{|\mathcal{I}_{b}|}$ the partial gradient w.r.t.\ that block.

At iteration $t=0,1,2,\dots$:
\begin{enumerate}[label=\arabic*)]
\item Sample a block $i_{t}\in\{1,\dots ,B\}$ uniformly at random,
      i.e. $\Pr(i_{t}=b)=\frac{1}{B}$.
\item Perform a gradient step on the selected block:
\[
w_{t+1}^{(b)}=
\begin{cases}
w_{t}^{(b)}-\eta_{t}\,\nabla_{b}f(w_{t}), & b=i_{t},\\[4pt]
w_{t}^{(b)}, & b\neq i_{t},
\end{cases}
\tag{1}
\]
where $\eta_{t}>0$ is a (possibly time‑varying) stepsize.
\end{enumerate}
All other blocks remain unchanged.  

\section*{Pseudocode}
\begin{algorithm}[H]
\caption{Randomized Block Coordinate Descent}
\begin{algorithmic}[1]
\State \textbf{Input:} initial point $w_{0}\in\R^{d}$, stepsize schedule $\{\eta_{t}\}_{t\ge0}$, number of blocks $B$.
\For{$t=0,1,\dots,T-1$}
    \State Sample $i_{t}\sim\text{Uniform}\{1,\dots ,B\}$.
    \State Compute partial gradient $g_{t}= \nabla_{i_{t}}f(w_{t})$.
    \State Update block $i_{t}$:
    \[
    w_{t+1}^{(i_{t})}= w_{t}^{(i_{t})}-\eta_{t}\,g_{t}.
    \]
    \State For all $b\neq i_{t}$ set $w_{t+1}^{(b)}=w_{t}^{(b)}$.
\EndFor
\State \textbf{Output:} $w_{T}$ (or the average $\bar w_{T}=\frac1T\sum_{t=0}^{T-1}w_{t}$).
\end{algorithmic}
\end{algorithm}

\section*{Dry Run Example}
Consider the scalar function $f(w)=(w-3)^{2}$, which we view as a single‑block problem ($B=1$).  
Take $w_{0}=0$, constant stepsize $\eta=0.1$, and run three iterations.

\begin{center}
\begin{tabular}{c|c|c|c}
$t$ & $w_{t}$ & $\nabla f(w_{t})$ & $f(w_{t})$\\\hline
0 & $0$ & $2(w_{0}-3)=-6$ & $9$\\
1 & $0-0.1(-6)=0.6$ & $2(0.6-3)=-4.8$ & $(0.6-3)^{2}=5.76$\\
2 & $0.6-0.1(-4.8)=1.08$ & $2(1.08-3)=-3.84$ & $(1.08-3)^{2}=3.6864$\\
3 & $1.08-0.1(-3.84)=1.464$ & $2(1.464-3)=-3.072$ & $(1.464-3)^{2}=2.3665$
\end{tabular}
\end{center}
The objective value decreases at each step, illustrating the algorithm’s progress.

\newpage
\section*{Assumptions}
\begin{assumption}[Convexity]\label{ass:conv}
$f$ is convex, i.e.
\[
f(y)\ge f(x)+\langle\nabla f(x),y-x\rangle,\qquad\forall x,y\in\R^{d}.
\]
\end{assumption}
\begin{assumption}[Block‑wise $L_{b}$‑smoothness]\label{ass:smooth}
For each block $b\in\{1,\dots ,B\}$ there exists $L_{b}>0$ such that for all $x,y$ differing only in block $b$,
\[
\|\nabla_{b}f(x)-\nabla_{b}f(y)\| \le L_{b}\,\|x^{(b)}-y^{(b)}\|.
\]
Equivalently,
\[
f(x+U_{b}h)\le f(x)+\langle\nabla_{b}f(x),h\rangle+\frac{L_{b}}{2}\|h\|^{2},
\tag{2}
\]
where $U_{b}$ embeds a block vector $h$ into $\R^{d}$ (zeros elsewhere).
\end{assumption}
\begin{assumption}[Bounded optimal distance]\label{ass:diam}
Let $w^{*}$ be a minimizer of $f$. Define a block‑wise diameter
\[
D^{2}:=\sum_{b=1}^{B}\|w^{*(b)}-w_{0}^{(b)}\|^{2}<\infty .
\]
\end{assumption}
\begin{assumption}[Bounded partial gradients]\label{ass:grad}
There exists $G>0$ such that for all $w$ and all blocks $b$,
\[
\|\nabla_{b}f(w)\|\le G .
\]
\end{assumption}

\section*{Main Convergence Theorem}
\begin{theorem}[Expected suboptimality of RBCD]\label{thm:main}
Assume \ref{ass:conv}–\ref{ass:grad} and use a *constant* stepsize
\[
\eta_{t}\equiv \eta \le \frac{1}{\max_{b}L_{b}} .
\]
Then for any $T\ge 1$ the iterate $w_{T}$ generated by (1) satisfies
\[
\E\!\left[\,f(w_{T})-f(w^{*})\,\right]
\le
\frac{2\,\sum_{b=1}^{B}L_{b}\,D^{2}}{B\,\eta\,T}
\;+\;
\frac{\eta\,G^{2}}{2},
\tag{3}
\]
where the expectation is taken over the random block selections $\{i_{t}\}_{t=0}^{T-1}$.
In particular, choosing
\[
\eta^{\star}= \sqrt{\frac{2\sum_{b}L_{b}\,D^{2}}{B\,G^{2}\,T}}
\]
balances the two terms and yields the optimal $O(1/\sqrt{T})$ rate:
\[
\E\!\left[\,f(w_{T})-f(w^{*})\,\right]\le
\sqrt{\frac{2\,\sum_{b}L_{b}\,D^{2}\,G^{2}}{B\,T}} .
\tag{4}
\]
\end{theorem}

\section*{Theoretical Properties}
\begin{itemize}
\item \textbf{Stability.} The stepsize restriction $\eta\le 1/\max_{b}L_{b}$ guarantees that the per‑block update (2) never increases the objective in expectation.
\item \textbf{Hyperparameter sensitivity.} The bound (3) depends on the *average* smoothness $\frac{1}{B}\sum_{b}L_{b}$ rather than the worst‑case $L_{\max}$, so the algorithm is less sensitive to a few ill‑conditioned blocks.
\item \textbf{Comparison to full‑gradient GD.} Full‑gradient GD with stepsize $1/L_{\max}$ yields $O(1/T)$ for smooth convex functions. The randomized block scheme attains a comparable $O(1/T)$ rate when the number of blocks $B$ is constant; for large $B$ the per‑iteration cost drops proportionally, giving a better overall *computational* complexity.
\item \textbf{Special case $B=1$.} The algorithm reduces to standard (full) gradient descent and (3) collapses to the classical bound $\frac{L D^{2}}{2\eta T}+\frac{\eta G^{2}}{2}$.
\end{itemize}

\section*{Discussion}
The theorem shows that, under mild convexity and block‑smoothness assumptions, randomized block coordinate descent converges in expectation at the same $O(1/\sqrt{T})$ rate as SGD for non‑strongly convex problems, while its per‑iteration work scales with the size of a single block. When $f$ is additionally $\mu$‑strongly convex, the same proof technique (with a stronger descent inequality) yields an exponential rate $O\!\big((1-\frac{\mu}{B\bar L})^{T}\big)$, where $\bar L:=\frac{1}{B}\sum_{b}L_{b}$.

\newpage
\section*{Preliminaries (Definitions and Lemmas)}
\begin{lemma}[Descent Lemma for a single block]\label{lem:descent}
Under Assumption \ref{ass:smooth}, for any $w\in\R^{d}$ and any block $b$,
\[
f\!\big(w-U_{b}\eta\nabla_{b}f(w)\big)
\le
f(w)-\eta\|\nabla_{b}f(w)\|^{2}
+\frac{L_{b}\eta^{2}}{2}\|\nabla_{b}f(w)\|^{2}.
\tag{5}
\]
\end{lemma}
\begin{proof}
Apply the smoothness inequality (2) with $h=-\eta\nabla_{b}f(w)$:
\[
f\!\big(w+U_{b}h\big)\le f(w)+\langle\nabla_{b}f(w),h\rangle+\frac{L_{b}}{2}\|h\|^{2}.
\]
Substituting $h=-\eta\nabla_{b}f(w)$ gives (5). ∎
\end{proof}

\begin{lemma}[Expected distance reduction]\label{lem:dist}
Let $w^{*}$ be a minimizer. For any $t\ge0$,
\[
\E_{i_{t}}\!\big[\|w_{t+1}-w^{*}\|^{2}\big]
=
\|w_{t}-w^{*}\|^{2}
-\frac{2\eta}{B}\,\langle\nabla f(w_{t}),w_{t}-w^{*}\rangle
+\frac{\eta^{2}}{B}\sum_{b=1}^{B}\|\nabla_{b}f(w_{t})\|^{2}.
\tag{6}
\]
\end{lemma}
\begin{proof}
Condition on $w_{t}$ and expand the squared norm using (1).  
Only the selected block $i_{t}$ changes:
\[
\|w_{t+1}-w^{*}\|^{2}
=
\|w_{t}-w^{*}\|^{2}
-2\eta\langle\nabla_{i_{t}}f(w_{t}),w_{t}^{(i_{t})}-w^{*(i_{t})}\rangle
+\eta^{2}\|\nabla_{i_{t}}f(w_{t})\|^{2}.
\]
Take expectation over $i_{t}$ (uniform distribution):
\[
\E_{i_{t}}[\cdot]
=
\|w_{t}-w^{*}\|^{2}
-\frac{2\eta}{B}\sum_{b=1}^{B}\langle\nabla_{b}f(w_{t}),w_{t}^{(b)}-w^{*(b)}\rangle
+\frac{\eta^{2}}{B}\sum_{b=1}^{B}\|\nabla_{b}f(w_{t})\|^{2}.
\]
Recognizing the inner product as $\langle\nabla f(w_{t}),w_{t}-w^{*}\rangle$ yields (6). ∎
\end{proof}

\section*{Main Proof (Step‑by‑Step Derivation)}
We prove Theorem \ref{thm:main} by chaining Lemmas \ref{lem:descent} and \ref{lem:dist}.

\begin{proof}
Let $\eta\le 1/\max_{b}L_{b}$ be fixed.  

\noindent\textbf{[Step 1]} Apply Lemma \ref{lem:descent} conditioned on the random block $i_{t}$ and then take expectation over $i_{t}$:
\[
\E_{i_{t}}\!\big[f(w_{t+1})\big]
\le
f(w_{t})
-\eta\,\E_{i_{t}}\!\big[\|\nabla_{i_{t}}f(w_{t})\|^{2}\big]
+\frac{\eta^{2}}{2}\E_{i_{t}}\!\big[L_{i_{t}}\|\nabla_{i_{t}}f(w_{t})\|^{2}\big].
\tag{7}
\]
Because $i_{t}$ is uniform,
\[
\E_{i_{t}}\!\big[\|\nabla_{i_{t}}f(w_{t})\|^{2}\big]
=
\frac{1}{B}\sum_{b=1}^{B}\|\nabla_{b}f(w_{t})\|^{2},
\quad
\E_{i_{t}}\!\big[L_{i_{t}}\|\nabla_{i_{t}}f(w_{t})\|^{2}\big]
=
\frac{1}{B}\sum_{b=1}^{B}L_{b}\|\nabla_{b}f(w_{t})\|^{2}.
\tag{8}
\]

\noindent\textbf{[Step 2]} Insert (8) into (7):
\[
\E_{i_{t}}\!\big[f(w_{t+1})\big]
\le
f(w_{t})
-\frac{\eta}{B}\sum_{b=1}^{B}\|\nabla_{b}f(w_{t})\|^{2}
+\frac{\eta^{2}}{2B}\sum_{b=1}^{B}L_{b}\|\nabla_{b}f(w_{t})\|^{2}.
\tag{9}
\]

\noindent\textbf{[Step 3]} Since $\eta\le 1/\max_{b}L_{b}$ we have $L_{b}\eta\le 1$ for all $b$, hence
\[
\frac{\eta^{2}}{2B}L_{b}\|\nabla_{b}f(w_{t})\|^{2}
\le
\frac{\eta}{2B}\|\nabla_{b}f(w_{t})\|^{2}.
\tag{10}
\]
Plugging (10) into (9) gives
\[
\E_{i_{t}}\!\big[f(w_{t+1})\big]
\le
f(w_{t})
-\frac{\eta}{2B}\sum_{b=1}^{B}\|\nabla_{b}f(w_{t})\|^{2}.
\tag{11}
\]

\noindent\textbf{[Step 4]} By convexity (Assumption \ref{ass:conv}) we have
\[
f(w_{t})-f(w^{*})
\le
\langle\nabla f(w_{t}),w_{t}-w^{*}\rangle .
\tag{12}
\]

\noindent\textbf{[Step 5]} Apply Lemma \ref{lem:dist} and take total expectation (over all randomness up to iteration $t$). Using (12) to replace the inner product:
\[
\E\!\big[\|w_{t+1}-w^{*}\|^{2}\big]
\le
\E\!\big[\|w_{t}-w^{*}\|^{2}\big]
- \frac{2\eta}{B}\,\E\!\big[f(w_{t})-f(w^{*})\big]
+\frac{\eta^{2}}{B}\E\!\big[\|\nabla f(w_{t})\|^{2}\big].
\tag{13}
\]

\noindent\textbf{[Step 6]} Bound the full gradient norm by the sum of block norms and then by $G$ (Assumption \ref{ass:grad}):
\[
\|\nabla f(w_{t})\|^{2}
=
\Big\|\sum_{b=1}^{B}\nabla_{b}f(w_{t})\Big\|^{2}
\le
B\sum_{b=1}^{B}\|\nabla_{b}f(w_{t})\|^{2}
\le
B\,G^{2}.
\tag{14}
\]

\noindent\textbf{[Step 7]} Substitute (14) into (13):
\[
\E\!\big[\|w_{t+1}-w^{*}\|^{2}\big]
\le
\E\!\big[\|w_{t}-w^{*}\|^{2}\big]
- \frac{2\eta}{B}\,\E\!\big[f(w_{t})-f(w^{*})\big]
+ \eta^{2} G^{2}.
\tag{15}
\]

\noindent\textbf{[Step 8]} Rearrange (15) to isolate the suboptimality term:
\[
\E\!\big[f(w_{t})-f(w^{*})\big]
\le
\frac{B}{2\eta}\Big(
\E\!\big[\|w_{t}-w^{*}\|^{2}\big]
-\E\!\big[\|w_{t+1}-w^{*}\|^{2}\big]\Big)
+\frac{B\eta G^{2}}{2}.
\tag{16}
\]

\noindent\textbf{[Step 9]} Sum (16) over $t=0,\dots ,T-1$ and telescope the distance terms:
\[
\sum_{t=0}^{T-1}\E\!\big[f(w_{t})-f(w^{*})\big]
\le
\frac{B}{2\eta}\Big(
\|w_{0}-w^{*}\|^{2}
-\E\!\big[\|w_{T}-w^{*}\|^{2}\big]\Big)
+ \frac{B\eta G^{2}T}{2}.
\tag{17}
\]

\noindent\textbf{[Step 10]} Drop the non‑negative final distance term and divide by $T$:
\[
\frac{1}{T}\sum_{t=0}^{T-1}\E\!\big[f(w_{t})-f(w^{*})\big]
\le
\frac{B\|w_{0}-w^{*}\|^{2}}{2\eta T}
+ \frac{\eta G^{2}}{2}.
\tag{18}
\]

\noindent\textbf{[Step 11]} Relate the initial distance to the block‑wise diameters defined in Assumption \ref{ass:diam}:
\[
\|w_{0}-w^{*}\|^{2}
=
\sum_{b=1}^{B}\|w_{0}^{(b)}-w^{*(b)}\|^{2}
\le D^{2}.
\tag{19}
\]

\noindent\textbf{[Step 12]} Insert (19) into (18) and multiply numerator and denominator by $\sum_{b}L_{b}$ (which is positive) to obtain the bound in the theorem:
\[
\E\!\big[f(\bar w_{T})-f(w^{*})\big]
\le
\frac{2\big(\sum_{b=1}^{B}L_{b}\big) D^{2}}{B\eta T}
+\frac{\eta G^{2}}{2},
\]
where $\bar w_{T}$ denotes the uniform average of the iterates (convexity permits replacing the average of function values by the function value at the average). This is exactly inequality \eqref{eq:3}.

\noindent\textbf{[Step 13]} To obtain the optimized rate, set the derivative of the RHS w.r.t.\ $\eta$ to zero, yielding
\[
\eta^{\star}= \sqrt{\frac{2\big(\sum_{b}L_{b}\big)D^{2}}{B G^{2}T}} .
\]
Plugging $\eta^{\star}$ back gives (4). ∎
\end{proof}

\section*{Rate Derivation (With Constants)}
From Theorem \ref{thm:main} and the optimal stepsize $\eta^{\star}$ we have
\[
\E\!\big[f(w_{T})-f(w^{*})\big]
\le
\underbrace{\sqrt{\frac{2\big(\sum_{b}L_{b}\big)D^{2}G^{2}}{B}}}_{\displaystyle C}
\; \frac{1}{\sqrt{T}}.
\]
Thus the convergence constant $C$ scales linearly with the square‑root of the average block smoothness $\big(\sum_{b}L_{b}\big)/B$ and with the initial distance $D$.

\section*{Special Cases}
\begin{enumerate}[label=\arabic*)]
\item \textbf{Single block $(B=1)$.} The algorithm coincides with gradient descent. The bound reduces to
\[
\E[f(w_{T})-f(w^{*})]\le \frac{2L D^{2}}{\eta T}+\frac{\eta G^{2}}{2},
\]
recovering the classic result.
\item \textbf{Uniform smoothness $L_{b}=L$ for all $b$.} Then $\sum_{b}L_{b}=BL$ and the rate becomes
\[
\E[f(w_{T})-f(w^{*})]\le
\frac{2L D^{2}}{\eta T}+\frac{\eta G^{2}}{2},
\]
identical to the single‑block case but with per‑iteration cost reduced by a factor $B$.
\item \textbf{Strong convexity.} If $f$ is $\mu$‑strongly convex, one can strengthen (12) to
\[
f(w_{t})-f(w^{*})\ge \frac{\mu}{2}\|w_{t}-w^{*}\|^{2},
\]
and repeating the steps above yields a linear rate
\[
\E\!\big[f(w_{T})-f(w^{*})\big]\le
\big(1-\tfrac{\mu\eta}{B}\big)^{T}\big(f(w_{0})-f(w^{*})\big).
\]
\end{enumerate}

\end{document}